\section{Results: heme oxidation enzyme (method foundation)}
\label{sec:res:heme}

\plan{\textbf{Claim.} Fine-tuned ProteinMPNN beats SSM and vanilla MPNN
re-sampling at finding active variants under matched experimental effort.
\textbf{Risk.} This is the foundational claim of the paper. If the comparison
is not clean, the entire methods narrative is in trouble --- this section needs
the cleanest, most defensible head-to-head. Status: \statusblocked.}

\subsection*{Block 1 --- Setup}
\begin{planlist}
  \item Backbone: $\langle$de novo heme oxidation enzyme, RFdiffusion
        origin$\rangle$.
  \item Starting design: $\langle$activity baseline$\rangle$.
  \item Three arms, matched library budget ($N$ variants ordered/tested):
        (i) SSM around starting design; (ii) vanilla MPNN re-sampling on
        backbone; (iii) fine-tuned MPNN (this work) + fragment library.
\end{planlist}
\todo{heme results --- setup paragraph}

\subsection*{Block 2 --- Headline figure}
\begin{planlist}
  \item Activity distribution across all three arms.
  \item Hit rate (fraction above starting design / above background).
  \item Best-variant comparison.
\end{planlist}
\todo{heme results --- main comparison figure + numbers}

\subsection*{Block 3 --- Mutation distance / sequence-space analysis}
\begin{planlist}
  \item Where in sequence space the best variants from each arm live.
  \item Demonstrate the SSM and vanilla-MPNN arms are confined to a near
        neighborhood while fine-tuned MPNN reaches farther.
\end{planlist}
\todo{heme results --- sequence-space analysis}

\subsection*{Block 4 --- Mechanism check (if available)}
\begin{planlist}
  \item Structural / spectroscopic / kinetic confirmation that the improvement
        is real catalysis, not artifact.
\end{planlist}
\todo{heme results --- mechanism check}
